\begin{frame}[t,fragile]{最適化問題}
  \begin{itemize}
    %\setlength{\itemsep}{1em}
  \item 最適化問題の種類
    \begin{itemize}
    \item 連続最適化問題 $\Leftarrow$ 目的関数が凸ではない場合、難しい
    \item 離散最適化(組み合せ最適化)問題 $\Leftarrow$ さらに難しい
    \end{itemize}
  \item 真の(大局的な)最小値(最大値)を求めるのは難しい(一般的には極値のみ)。多次元では極小(大)を囲い込むことさえできない
  \item 導関数を使う方法: ニュートン法、準ニュートン法、最急降下法、勾配降下法、共役勾配法$\cdots$
  \item 使わない方法: 囲い込み法、Nelder-Meadの滑降シンプレックス法、シミュレーテッドアニーリング、量子アニーリング$\cdots$
  \item 目的関数・導関数の評価回数と収束までの反復回数のトレードオフ \\[1em]
  \item 「いかに優れた部分最適も全体最適には勝てない」-- P. F. Drucker
  \item 「あらゆる問題を効率よく解く万能アルゴリズムは存在しない (No-Free-Lunch Theorem)」-- D. H. Wolpert and W. G. Macready

  \end{itemize}
\end{frame}
